\section{Mathematical model}
\label{sec:model}

We describe the model in normalized form. The four-body central-configuration
formulation of Section~\ref{subsec:cc} is due to \cite{shoaib2023}
and is \emph{not} claimed as a result of the present paper; we restate it only
to fix notation and because it defines the parameter family along which the
restricted problem is posed. The restricted fifth-body model follows in
Sections~\ref{subsec:potential}--\ref{subsec:jacobi}.

\subsection{The primary central configuration}
\label{subsec:cc}

Following the normalization of Shoaib et al.\ \cite{shoaib2023}, we use the
translation, rotation, reflection, and scaling invariances of the planar
central-configuration problem to place the four primaries at
\begin{equation}
\mathbf r_0=(0,b),\qquad
\mathbf r_1=(-1,0),\qquad
\mathbf r_2=(0,a),\qquad
\mathbf r_3=(c_1,c_2),
\label{eq:primary_positions}
\end{equation}
with
\[
a>0,\qquad b<a,\qquad b\neq \pm a.
\]
The choice $\mathbf r_1=(-1,0)$ fixes the length scale, while
$\mathbf r_0$ and $\mathbf r_2$ lie on the $y$-axis and
$\mathbf r_3=(c_1,c_2)$ remains free in the plane. The geometry is therefore
described by the parameter vector
\[
\mathbf p=(a,b,c_1,c_2).
\]

For compactness, define
\begin{equation}
\begin{aligned}
\sigma &= (1+a^2)^{3/2},
&
\varsigma &= (1+b^2)^{3/2},
&
\tau &= (a-b)^3,
\\
G &= \left[c_1^2+(c_2-a)^2\right]^{3/2},
&
T &= \left[c_1^2+(c_2-b)^2\right]^{3/2},
&
M &= \left[(c_1+1)^2+c_2^2\right]^{3/2}.
\end{aligned}
\label{eq:auxiliary_distances}
\end{equation}
These quantities are the cubes of the six mutual distances between the
primaries.

Let
\[
\rho_0=\frac{m_0}{m_3},\qquad
\rho_1=\frac{m_1}{m_3},\qquad
\rho_2=\frac{m_2}{m_3}
\]
denote the primary mass ratios. For the nondegenerate geometries considered
here, the coplanar four-body central-configuration condition of
\cite{shoaib2023} reduces to the scalar equation
\begin{equation}
\begin{aligned}
\psi(a,b,c_1,c_2)
={}&
\tau(\sigma-\varsigma)M
+T\left[(\tau-\sigma)M+\sigma(\varsigma-\tau)\right]
\\
&+
G\left[(\sigma-\varsigma)T+(\varsigma-\tau)M
+\varsigma(\tau-\sigma)\right]
=0,
\end{aligned}
\label{eq:cc_constraint}
\end{equation}
and the corresponding mass ratios are
\begin{equation}
\rho_0=
\frac{(c_2-a-ac_1)(M-G)\varsigma\tau}
     {(a-b)(\varsigma-\tau)MG},
\qquad
\rho_1=
\frac{c_1(G-T)\sigma\varsigma}
     {(\sigma-\varsigma)GT},
\qquad
\rho_2=
\frac{(c_2-b-bc_1)(M-T)\sigma\tau}
     {(a-b)(\tau-\sigma)MT}.
\label{eq:mass_ratios}
\end{equation}
Equations~\eqref{eq:cc_constraint}--\eqref{eq:mass_ratios} are the
four-body central-configuration result of \cite{shoaib2023}; they are
restated here only to define the family of primary configurations used in
the restricted problem.

Not every solution of $\psi=0$ corresponds to a positive-mass,
nondegenerate configuration. We therefore define the admissible set
\begin{equation}
\mathcal A=
\left\{
\mathbf p:
\psi(\mathbf p)=0,\;
\rho_0(\mathbf p)>0,\;
\rho_1(\mathbf p)>0,\;
\rho_2(\mathbf p)>0
\right\},
\label{eq:admissible_set}
\end{equation}
subject additionally to non-vanishing denominators in
\eqref{eq:mass_ratios} and the exclusion of primary collisions. Since only
mass ratios enter the central-configuration construction, the overall mass
scale remains arbitrary. We fix this scale by setting $m_3=1$, so that the
normalized primary masses used throughout the restricted problem are
\[
(m_0,m_1,m_2,m_3)=(\rho_0,\rho_1,\rho_2,1).
\]

\subsection{The restricted fifth body: potential, equilibria, and linear stability}
\label{subsec:potential}

For each admissible primary configuration
$\mathbf p\in\mathcal A$, the four primary masses and positions are fixed in
the uniformly rotating frame, and the fifth body is taken to be
infinitesimal: it moves in their gravitational field but does not perturb
the primary central configuration. Let $(x,y)$ denote its rotating-frame
coordinates, and define the distances to the four primaries by
\begin{equation}
\begin{aligned}
r_{40} &= \sqrt{x^2+(y-b)^2},
&
r_{41} &= \sqrt{(x+1)^2+y^2},
\\
r_{42} &= \sqrt{x^2+(y-a)^2},
&
r_{43} &= \sqrt{(x-c_1)^2+(y-c_2)^2}.
\end{aligned}
\label{eq:test_particle_distances}
\end{equation}

With the rotation rate normalized to unity, the effective potential is
\begin{equation}
\Omega(x,y)
=
\frac{1}{2}(x^2+y^2)
+
\sum_{i=0}^{3}\frac{m_i}{r_{4i}},
\label{eq:effective_potential}
\end{equation}
where
\[
(m_0,m_1,m_2,m_3)
=
(\rho_0,\rho_1,\rho_2,1).
\]
The equations of motion are
\begin{equation}
\ddot{x}-2\dot{y}=\Omega_x,
\qquad
\ddot{y}+2\dot{x}=\Omega_y,
\label{eq:equations_of_motion}
\end{equation}
with
\begin{equation}
\Omega_x
=
x-\sum_{i=0}^{3}m_i\frac{x-x_i}{r_{4i}^3},
\qquad
\Omega_y
=
y-\sum_{i=0}^{3}m_i\frac{y-y_i}{r_{4i}^3},
\label{eq:potential_gradient}
\end{equation}
where $(x_i,y_i)$ are the primary positions defined in
\eqref{eq:primary_positions}.

The relative equilibria, or libration points, of the infinitesimal body are
the stationary solutions of the rotating-frame equations,
\begin{equation}
\Omega_x(x,y)=0,
\qquad
\Omega_y(x,y)=0.
\label{eq:equilibrium_equations}
\end{equation}
Unlike the five Lagrange points of the CR3BP, the number and locations of
these equilibria depend on the primary configuration
$\mathbf p\in\mathcal A$. Their variation along the admissible
central-configuration set is the equilibrium bifurcation problem studied in
Section~\ref{sec:results}.

To analyse both the local root structure of
\eqref{eq:equilibrium_equations} and the linear dynamics near an
equilibrium, we use the Hessian of $\Omega$. Writing
\[
\Delta x_i=x-x_i,
\qquad
\Delta y_i=y-y_i,
\]
its entries are
\begin{equation}
\begin{aligned}
\Omega_{xx}
&=
1-\sum_{i=0}^{3}m_i
\left(
\frac{1}{r_{4i}^3}
-
\frac{3\Delta x_i^2}{r_{4i}^5}
\right),
\\
\Omega_{yy}
&=
1-\sum_{i=0}^{3}m_i
\left(
\frac{1}{r_{4i}^3}
-
\frac{3\Delta y_i^2}{r_{4i}^5}
\right),
\\
\Omega_{xy}
&=
3\sum_{i=0}^{3}m_i
\frac{\Delta x_i\Delta y_i}{r_{4i}^5}.
\end{aligned}
\label{eq:hessian_entries}
\end{equation}
Thus
\begin{equation}
H_\Omega
=
\begin{pmatrix}
\Omega_{xx} & \Omega_{xy}\\
\Omega_{xy} & \Omega_{yy}
\end{pmatrix}
\label{eq:hessian}
\end{equation}
is simultaneously the Hessian of the effective potential and the Jacobian
of the equilibrium map
\[
F(x,y)=
\begin{pmatrix}
\Omega_x\\
\Omega_y
\end{pmatrix}.
\]
Consequently, $\det H_\Omega=0$ marks a degeneracy of the equilibrium
equations and plays a central role in the fold analysis of Sections~\ref{sec:methods}--\ref{sec:results}.

Let $(x_e,y_e)$ be an equilibrium. Linearizing
\eqref{eq:equations_of_motion} about $(x_e,y_e)$ gives
\begin{equation}
\dot{\mathbf X}=A\mathbf X,
\qquad
\mathbf X=
\begin{pmatrix}
X\\Y\\\dot X\\\dot Y
\end{pmatrix},
\end{equation}
with
\begin{equation}
A=
\begin{pmatrix}
0 & 0 & 1 & 0\\
0 & 0 & 0 & 1\\
\Omega_{xx} & \Omega_{xy} & 0 & 2\\
\Omega_{xy} & \Omega_{yy} & -2 & 0
\end{pmatrix}_{(x_e,y_e)} .
\label{eq:variational_matrix}
\end{equation}
Its characteristic polynomial is
\begin{equation}
\lambda^4
+
\left(4-\Omega_{xx}-\Omega_{yy}\right)\lambda^2
+
\left(\Omega_{xx}\Omega_{yy}-\Omega_{xy}^2\right)
=0.
\label{eq:characteristic_polynomial}
\end{equation}
Because \eqref{eq:characteristic_polynomial} is an even polynomial with real
coefficients, the spectrum is symmetric under both
$\lambda\mapsto-\lambda$ and complex conjugation. Hence the four
eigenvalues occur as
\[
\{\lambda,-\lambda,\bar{\lambda},-\bar{\lambda}\},
\]
with the usual reductions when some members coincide.

An equilibrium is spectrally unstable whenever at least one eigenvalue has
positive real part. It is spectrally stable when the spectrum consists of
two purely imaginary conjugate pairs,
\[
\lambda=\pm i\omega_1,
\qquad
\lambda=\pm i\omega_2,
\qquad
\omega_1,\omega_2>0.
\]
We use ``spectrally stable'' and ``linearly stable'' in this restricted
sense throughout. This condition concerns only the linearized dynamics and
does not by itself imply nonlinear or KAM stability. The centre frequencies
$\omega_1$ and $\omega_2$ identified here provide the linear modes from
which the periodic-orbit families of Section~\ref{sec:periodic} are initialized.


\subsection{Jacobi integral and Hill regions}
\label{subsec:jacobi}

The rotating-frame equations \eqref{eq:equations_of_motion} admit the
Jacobi integral. Multiplying the first equation by $\dot{x}$ and the second
by $\dot{y}$, and then adding, gives
\[
\dot{x}\ddot{x}+\dot{y}\ddot{y}
=
\Omega_x\dot{x}+\Omega_y\dot{y},
\]
because the Coriolis terms cancel identically. Hence
\[
\frac{d}{dt}
\left[
\frac{1}{2}\left(\dot{x}^2+\dot{y}^2\right)-\Omega(x,y)
\right]
=0.
\]
We therefore adopt the standard Jacobi constant
\begin{equation}
C_J
=
2\Omega(x,y)
-
\left(\dot{x}^2+\dot{y}^2\right),
\label{eq:jacobi_integral}
\end{equation}
which is constant along every trajectory.

Since the kinetic term is non-negative, motion at a prescribed value of
$C_J$ is possible only where
\begin{equation}
2\Omega(x,y)\geq C_J.
\label{eq:hill_condition}
\end{equation}
The corresponding Hill region is
\begin{equation}
\mathcal H(C_J)
=
\left\{
(x,y):2\Omega(x,y)\geq C_J
\right\},
\label{eq:hill_region}
\end{equation}
and its boundary,
\begin{equation}
2\Omega(x,y)=C_J,
\label{eq:zero_velocity_curve}
\end{equation}
is the zero-velocity curve.

At an equilibrium $(x_e,y_e)$ the rotating-frame velocity vanishes, so its
critical Jacobi value is
\begin{equation}
C_{J,e}
=
2\Omega(x_e,y_e).
\label{eq:critical_jacobi}
\end{equation}
These equilibrium values provide the critical levels at which the geometry
of the regular Hill region can change through the opening or closing of
necks. Their relation to the equilibrium branches and to a representative
Hill-region topology change is examined in Section~\ref{sec:hill}.